\section*{Conclusion}

We present the first disk bispectrum inversion, which reveals the minimal sufficient coefficients needed for accurate shape recovery. This enables the construction of our \emph{selective} disk bispectrum, which is \textbf{17,322× faster} to compute than the full bispectrum for $(122\times112)$ images. Crucially, our disk bispectrum inversion provides a mapping from disk bispectrum coefficients back to the image they represent, meaning that statistics can be done in bispectrum space and the result can be mapped back to image space. Without this, methods such as MRA would be impossible. Accordingly, we additionally derive the noise propagation that would occur for selective disk bispectrum MRA (in line with \cite{bendory2018bispectrum} and use it to showcase our bispectrum inverse in MRA experiments. The selective disk bispectrum's speed and invertibility position it as a tractable and scalable shape representation for rotation-invariant learning.
